\documentclass[11pt]{article}
\usepackage[utf8]{inputenc}
\usepackage[T1]{fontenc}
\usepackage{amsmath,booktabs,graphicx}
\usepackage[margin=1in]{geometry}
\usepackage{hyperref}

\title{A Reproducible Computational Research Example}
\author{First Author \and Second Author}
\date{}

\begin{document}
\maketitle

\begin{abstract}
This example illustrates a reusable scientific article structure in LaTeX.
\end{abstract}

\section{Introduction}
Computational workflows can improve research transparency and reproducibility \cite{wilson2014best}.

\section{Methods}
Define observed change as
\begin{equation}
\Delta_i=y_i^{(\mathrm{after})}-y_i^{(\mathrm{before})}.
\end{equation}

\section{Results}

\begin{table}[ht]
\centering
\caption{Illustrative results by group.}
\begin{tabular}{lcc}\toprule
Group & Sample size & Mean change\\\midrule
A&42&2.47\\ B&42&3.67\\ C&42&4.60\\\bottomrule
\end{tabular}
\end{table}

\begin{figure}[ht]
\centering
\fbox{\parbox[c][1.2in][c]{0.78\textwidth}{\centering
Research Question $\rightarrow$ Data $\rightarrow$ Analysis $\rightarrow$ Evidence $\rightarrow$ Communication}}
\caption{Simplified reproducible research workflow.}
\end{figure}

\section{Discussion}
A documented path from question to evidence supports reproducibility.

\section{Conclusion}
LaTeX helps maintain consistent scientific documents as content evolves.

\bibliographystyle{plain}
\bibliography{references}

\end{document}
